\documentclass[
  digital, %% The `digital` option enables the default options for the
           %% digital version of a document. Replace with `printed`
           %% to enable the default options for the printed version
           %% of a document.
%%  color,   %% Uncomment these lines (by removing the %% at the
%%           %% beginning) to use color in the printed version of your
%%           %% document
  oneside, %% The `oneside` option enables one-sided typesetting,
           %% which is preferred if you are only going to submit a
           %% digital version of your thesis. Replace with `twoside`
           %% for double-sided typesetting if you are planning to
           %% also print your thesis. For double-sided typesetting,
           %% use at least 120 g/m² paper to prevent show-through.
  lof,     %% The `lof` option prints the List of Figures. Replace
           %% with `nolof` to hide the List of Figures.
  nolot,     %% The `lot` option prints the List of Tables. Replace
           %% with `nolot` to hide the List of Tables.
]{fithesis4}
%% The following section sets up the locales used in the thesis.
\usepackage[resetfonts]{cmap} %% We need to load the T2A font encoding
\usepackage[T1,T2A]{fontenc}  %% to use the Cyrillic fonts with Russian texts.
\usepackage[
  main=english, %% By using `czech` or `slovak` as the main locale
                %% instead of `english`, you can typeset the thesis
                %% in either Czech or Slovak, respectively.
  english, german, russian, czech, slovak %% The additional keys allow
]{babel}        %% foreign texts to be typeset as follows:
%%
%%   \begin{otherlanguage}{german}  ... \end{otherlanguage}
%%   \begin{otherlanguage}{russian} ... \end{otherlanguage}
%%   \begin{otherlanguage}{czech}   ... \end{otherlanguage}
%%   \begin{otherlanguage}{slovak}  ... \end{otherlanguage}
%%
%% For non-Latin scripts, it may be necessary to load additional
%% fonts:
\usepackage{paratype}
\def\textrussian#1{{\usefont{T2A}{PTSerif-TLF}{m}{rm}#1}}
%%
%% The following section sets up the metadata of the thesis.
\thesissetup{
    date        = \the\year/\the\month/\the\day,
    university  = mu,
    faculty     = fi,
    type        = mgr,
    department  = Department of Computer Systems and Communications,
    author      = Alexander Zgabur,
    gender      = m,
    advisor     = {RNDr. Adam Rambousek, Ph.D.},
    title       = {Generation and Automation of a Testing Matrix for Kuadrant project},
    TeXtitle    = {Generation and Automation of a Testing Matrix for Kuadrant project},
    keywords    = {Tekton, CI/CD Kubernetes},
    TeXkeywords = {Tekton, CI/CD, Kubernetes},
    abstract    = {%
        Todo
    },
    thanks      = {%
        Ing. Zdeněk Kraus a spol.
    },
    bib         = thesis.bib,
    %% Remove the following line to use the JVS 2018 faculty logo.
    facultyLogo = fithesis-fi,
}
\usepackage{makeidx}      %% The `makeidx` package contains
\makeindex                %% helper commands for index typesetting.
\usepackage[acronym]{glossaries}          %% The `glossaries` package
\renewcommand*\glspostdescription{\hfill} %% contains helper commands
\loadglsentries{thesis-terms-abbrs.tex}  %% for typesetting glossaries
\makenoidxglossaries                      %% and lists of abbreviations.
%% These additional packages are used within the document:
\usepackage{paralist} %% Compact list environments
\usepackage{amsmath}  %% Mathematics
\usepackage{amsthm}
\usepackage{amsfonts}
\usepackage{url}      %% Hyperlinks
\usepackage{markdown} %% Lightweight markup
\usepackage{listings} %% Source code highlighting
\lstset{
  basicstyle      = \ttfamily,
  identifierstyle = \color{black},
  keywordstyle    = \color{blue},
  keywordstyle    = {[2]\color{cyan}},
  keywordstyle    = {[3]\color{olive}},
  stringstyle     = \color{teal},
  commentstyle    = \itshape\color{magenta},
  breaklines      = true,
}
\usepackage{floatrow} %% Putting captions above tables
\floatsetup[table]{capposition=top}
\usepackage[babel]{csquotes} %% Context-sensitive quotation marks
\usepackage{float} %% for table, figure positioning
\usepackage{svg} %% for svg figures
\usepackage{enumerate} %% for better enumeration lists
\begin{document}
%% Uncomment the following lines (by removing the %% at the beginning)
%% and to print out List of Abbreviations and/or Glossary in your
%% document. Titles for these tables can be changed by replacing the
%% titles `Abbreviations` and `Glossary`, respectively.
\clearpage
\printnoidxglossary[title={Abbreviations}, type=\acronymtype]
%% \printnoidxglossary[title={Glossary}] % not used yet

\chapter*{Introduction}
\addcontentsline{toc}{chapter}{Introduction}
%Intro \acrfull{DNS}

\chapter{Background and goal}
\label{ch:first}
Intro to chapter
\section{Automated matrix testing}
What it is, and the benefits
\section{Kuadrant project}
What is Kuadrant
\subsection{Kuadrant supported matrix}
The matrix

\chapter{Pairwise testing}
Intro to chapter.
\section{Pairwise combination theory}
What is pairwise, why is it a fit to Kuadrant.
\section{Analysis of existing solution (jurban work)}

\chapter{Improvements to generator}
1. Weights do not work. Prioritization of specific configuration should be improved.
2. Mandatory pairs, Exluded pairs, max occurances of confuguration as opposed to weights
3. json input -> yaml input
4. Tekton plugin which would run pipeline(s) for achieving the testing of configurations
5. Architecture support
6. Dockerize
7. Tests + github actions

\chapter{Test automation}
Into into chapter
\section{CI/CD pipelines in general}

\section{Tekton}
\section{Integration of generator and Tekton}

\chapter{Evaluation}
Example of generated configurations and tekton pipeline runs for single release of Kuadrant.

\chapter{Future work}
GUI

\chapter*{Conclusion}
\addcontentsline{toc}{chapter}{Conclusion}
Lorem ipsum


\appendix
\chapter{Electronic appendix}
\label{ch:e_appendix}
Electronic appendix of this thesis contains:
\begin{itemize}
    \item ...
\end{itemize}

\end{document}
